\subsection{p-FILTRATION OF FREE GROUPS}

In this no., \(p\) denotes a prime number and we assume that \(\mathbf{K} = \mathbf{F}_p\) . Let \(\gamma\) be the homomorphism of \(\mathrm{F}(\mathbf{X})\) into \(\Gamma(\mathbf{X})\) defined by \(\gamma(x) = 1 + x\) for \(x\) in \(\mathbf{X}\) ; we write \(\mathrm{F}_n^{(p)}(\mathbf{X}) = \gamma(1 + \hat{\Lambda}_n(\mathbf{X}))\) . The sequence \((\mathrm{F}_n^{(p)}(\mathbf{X}))_{n\geq 1}\) is an integral central filtration on \(\mathrm{F}(\mathbf{X})\) , which is separated since \(\gamma\) is injective (no. 3, Theorem 1). It is called the \(p\) -filtration on \(\mathrm{F}(\mathbf{X})\) .

PROPOSITION 2. Suppose that X is finite. For every integer \(n \geqslant 1\) , the group \(F(X) / F_n^{(p)}(X)\) is a finite \(p\) -group of nilpotency class \(\leqslant n\) .

Arguing by induction on \(n\) , it suffices to prove that \(\mathrm{F}_{n}^{(p)}(\mathbf{X}) / \mathrm{F}_{n + 1}^{(p)}(\mathbf{X})\) is a finite commutative \(p\) -group for all \(n \geqslant 1\) . For all \(w \in \mathrm{F}_{n}^{(p)}(\mathbf{X})\) , the element \(\gamma(w) - 1\) of \(\hat{\mathbf{A}}(\mathbf{X})\) is of order \(\geqslant n\) ; we denote by \(\delta_n(w)\) the homogeneous component of \(\gamma(w) - 1\) of degree \(n\) . The mapping \(\delta_n: \mathrm{F}_{n}^{(p)}(\mathbf{X}) \to \mathrm{A}^n (\mathbf{X})\) is a homomorphism with kernel \(\mathrm{F}_{n + 1}^{(p)}(\mathbf{X})\) (§ 4, no. 5, Proposition 3) and hence \(\mathrm{F}_{n}^{(p)}(\mathbf{X}) / \mathrm{F}_{n + 1}^{(p)}(\mathbf{X})\) is isomorphic to a subgroup of \(\mathrm{A}^n (\mathbf{X})\) . Since \(\mathbf{X}\) is finite, \(\mathrm{A}^n (\mathbf{X})\) is a finite-dimensional vector space over \(\mathbf{F}_p\) and hence a finite commutative \(p\) -group and so is \(\mathrm{F}_{n}^{(p)}(\mathbf{X}) / \mathrm{F}_{n + 1}^{(p)}(\mathbf{X})\) .

PROPOSITION 3. For all \(w \neq 1\) in \(\mathrm{F}(\mathbf{X})\) , there exist a finite \(p\) -group \(\mathbf{G}\) and a homomorphism \(f\) of \(\mathrm{F}(\mathbf{X})\) into \(\mathbf{G}\) such that \(f(w) \neq 1\) .

There exist elements \(x_{1},\ldots ,x_{r}\) of \(\mathbf{X}\) and integers \(n_1,\dots ,n_r\) such that \(w = x_1^{n_1}\dots x_r^{n_r}\) . Let \(\mathrm{Y} = \{x_1,\dots,x_r\}\) . The canonical injection of \(\mathrm{Y}\) into \(\mathbf{X}\) extends to a homomorphism \(\alpha : \mathrm{F}(\mathrm{Y})\to \mathrm{F}(\mathrm{X})\) ; on the other hand, let \(\beta\) be the homomorphism of \(\mathrm{F}(\mathrm{X})\) into \(\mathrm{F}(\mathrm{Y})\) whose restriction to \(\mathbf{Y}\) is the identity and which maps \(\mathbf{X} - \mathbf{Y}\) to \(\{1\}\) . Then \(\beta (\alpha (y)) = y\) for \(y\in \mathbf{Y}\) and hence \(\beta \circ \alpha\) is the identity automorphism of \(\mathrm{F}(\mathrm{Y})\) . There obviously exists \(w^{\prime}\) in \(\mathrm{F}(\mathrm{Y})\) such that \(w = \alpha (w^{\prime})\) ; then \(\beta (w) = w^{\prime}\neq 1\) ; now \(\bigcap_{n\geqslant 1}\mathbf{F}_n^{(p)}(\mathbf{Y}) = \{1\}\) and there therefore exists an integer \(n\geqslant 1\) such that \(\beta (w)\notin \mathbf{F}_{\mathrm{n}}^{(p)}(\mathbf{Y})\) . By Proposition 2, the group \(\mathbf{G} = \mathbf{F}(\mathbf{Y}) / \mathbf{F}_{\mathrm{n}}^{(p)}(\mathbf{Y})\) is a finite \(p\) -group. If \(f\) is the composition of \(\beta\) and the canonical homomorphism of \(\mathrm{F}(\mathrm{Y})\) onto \(\mathbf{G}\) , then \(f(w)\neq 1\) .

COROLLARY. The intersection of the normal subgroups of finite index in F(X) is \{1\}.

